% =============================================================================
% RAPPORT PROFESSIONNEL — UrbanOps
% Thème : rouge / vert | Cadres : fond blanc, contours uniquement
% Compilation : pdflatex x2 (depuis docs/latex/)
% =============================================================================
% openany : évite les pages blanches volontaires avant chaque chapitre (openright les imposait)
\documentclass[12pt,a4paper,twoside,openany]{book}

\usepackage[T1]{fontenc}
\usepackage[utf8]{inputenc}
\usepackage[french]{babel}
\usepackage{microtype}

\usepackage[a4paper,left=2.5cm,right=2.5cm,top=2.5cm,bottom=2.5cm,bindingoffset=0.5cm]{geometry}
\usepackage{setspace}
\setstretch{1.2}
\usepackage{ragged2e}
\usepackage{enumitem}
\setlist{itemsep=0.32em, parsep=0pt, topsep=0.35em}

\usepackage{graphicx}
\usepackage[table]{xcolor}
% --- Palette demandée (rouge + vert) ---
\definecolor{UrbanRed}{HTML}{B91C1C}
\definecolor{UrbanGreen}{HTML}{15803D}
\definecolor{UrbanDark}{HTML}{1F2937}
\definecolor{FrameNeutral}{HTML}{374151}
\definecolor{CodeBg}{HTML}{FAFAFA}

\usepackage{float}
\usepackage{placeins}
\usepackage{caption}
\captionsetup{font=small,labelfont={bf,color=UrbanRed}}
\captionsetup[table]{labelfont={bf,color=UrbanGreen}}

\usepackage{booktabs}
\usepackage{array}
\usepackage{longtable}
\usepackage{tabularx}
\newcolumntype{Y}{>{\RaggedRight\arraybackslash}X}

\usepackage{amsmath,amssymb}

\usepackage{listings}
\lstdefinestyle{urban}{%
  backgroundcolor=\color{white},
  basicstyle=\ttfamily\footnotesize,
  breaklines=true,
  frame=single,
  framerule=0.45pt,
  rulecolor=\color{FrameNeutral},
  columns=fullflexible,
  keepspaces=true,
  showstringspaces=false,
  tabsize=2,
}
\lstset{style=urban}
\lstdefinestyle{sqlurban}{style=urban,language=SQL}
\lstdefinestyle{javaurban}{style=urban,language=Java}
\lstdefinestyle{bashurban}{style=urban,language=bash}

\usepackage[most]{tcolorbox}
\tcbuselibrary{skins,breakable}

\usepackage{tocloft}
\renewcommand{\cftchapfont}{\bfseries\color{UrbanRed}}
\renewcommand{\cftsecfont}{\color{UrbanDark}}
\renewcommand{\cftsubsecfont}{\color{UrbanDark!85}}
\setlength{\cftbeforechapskip}{0.5em}
\renewcommand{\cftfigfont}{\color{UrbanDark}}
\renewcommand{\cfttabfont}{\color{UrbanDark}}

\usepackage{titlesec}
\titleformat{\chapter}[display]
  {\normalfont\huge\bfseries\color{UrbanRed}}
  {\chaptertitlename~\thechapter}
  {10pt}
  {\Huge}
\titlespacing*{\chapter}{0pt}{-6pt}{20pt}
\titleformat{name=\chapter,numberless}[display]
  {\normalfont\huge\bfseries\color{UrbanRed}}
  {}
  {0pt}
  {\Huge}
\titlespacing*{name=\chapter,numberless}{0pt}{-6pt}{20pt}
\titleformat{\section}
  {\Large\bfseries\color{UrbanGreen}}
  {\thesection}{0.7em}{}
\titlespacing*{\section}{0pt}{1em}{0.55em}
\titleformat{\subsection}
  {\large\bfseries\color{UrbanDark}}
  {\thesubsection}{0.6em}{}

\setcounter{tocdepth}{3}
\setcounter{secnumdepth}{3}

\usepackage[hidelinks,pdfusetitle]{hyperref}
\hypersetup{
  colorlinks=true,
  linkcolor=UrbanRed!90!black,
  urlcolor=UrbanGreen!80!black,
  citecolor=UrbanRed,
  bookmarksnumbered=true,
  pdftitle={UrbanOps - Rapport professionnel},
  pdfauthor={Equipe UrbanOps},
}
\usepackage{bookmark}
\usepackage[nameinlink]{cleveref}

\usepackage{fancyhdr}
\pagestyle{fancy}
\fancyhf{}
\fancyhead[LE]{\small\textcolor{UrbanDark}{UrbanOps — supervision urbaine}}
\fancyhead[RO]{\small\textcolor{UrbanDark}{2025--2026}}
\fancyhead[LO,RE]{\small\leftmark}
\fancyfoot[C]{\thepage}
\renewcommand{\headrulewidth}{0.4pt}
\renewcommand{\headrule}{\textcolor{UrbanGreen}{\rule{\headwidth}{\headrulewidth}}}
\renewcommand{\footrulewidth}{0pt}
\makeatletter
\let\ps@plain\ps@fancy
\makeatother

% =============================================================================
% Encadrés : fond BLANC, bordure colorée uniquement (pas de fond teinté)
% =============================================================================
\newtcolorbox{resumebox}{
  enhanced,
  colback=white,
  colframe=UrbanRed,
  boxrule=0.55pt,
  arc=2mm,
  left=7mm,right=7mm,top=5mm,bottom=5mm,
  breakable,
}

\newtcolorbox{archibox}[1][]{
  enhanced,
  colback=white,
  colframe=UrbanGreen,
  fonttitle=\bfseries\color{UrbanRed},
  title={#1},
  attach boxed title to top left={xshift=6mm,yshift=-2.5mm},
  boxed title style={colback=white,colframe=UrbanGreen,boxrule=0.55pt},
  arc=2mm,
  boxrule=0.55pt,
  breakable,
}

\newtcolorbox{notebox}{
  colback=white,
  colframe=UrbanRed!70!black,
  boxrule=0.45pt,
  arc=1.5mm,
  left=5mm,right=5mm,top=4mm,bottom=4mm,
  breakable,
}

% Inclusion figure OU cadre réservé (fond blanc)
\newcommand{\urbanfig}[4]{%
  \begin{figure}[!htbp]%
    \centering%
    \IfFileExists{#2}{%
      \includegraphics[width=#3\textwidth]{#2}%
    }{%
      \begin{tcolorbox}[
        enhanced,colback=white,colframe=UrbanGreen,
        width=0.92\textwidth,arc=2mm,halign=center,valign=center,
        height=0.28\textheight,boxrule=0.55pt]
        {\large\bfseries\color{UrbanRed} Capture à insérer}\\[2mm]
        {\small\texttt{#2}}\\[1.5mm]
        {\footnotesize Exporter votre capture sous ce nom exact dans \texttt{docs/latex/figures/}.}\\[1mm]
        {\small\textit{Légende prévue :} #4}
      \end{tcolorbox}%
    }%
    \caption{#4}%
    \label{#1}%
  \end{figure}%
}

\newcounter{prechapter}
\setcounter{prechapter}{0}
\newcommand{\ChapterRoman}[1]{%
  \refstepcounter{prechapter}%
  \phantomsection
  \chapter*{Chapitre \Roman{prechapter} — #1}%
  \addcontentsline{toc}{chapter}{Chapitre \Roman{prechapter} — #1}%
  \markboth{Chapitre \Roman{prechapter} — #1}{}%
}

\newcommand{\SecRoman}[1]{%
  \phantomsection
  \section*{#1}%
  \addcontentsline{toc}{section}{#1}%
}

\title{Rapport de projet — UrbanOps\\\large Plateforme de supervision et de signalement d'incidents urbains}
\author{\'Equipe UrbanOps}
\date{\today}

\begin{document}

\begin{titlepage}
  \thispagestyle{empty}
  \newgeometry{margin=2.2cm}
  \centering
  \null\vspace*{0.6cm}
  {\color{UrbanGreen}\rule{0.78\textwidth}{1.8pt}}\par\vspace{0.5cm}
  {\Large\bfseries\textsc{Institution / \'Ecole}}\par\vspace{0.15cm}
  {\large Département — filière / module}\par\vspace{0.3cm}
  {\normalsize\itshape Rapport technique — projet intégrateur / PFE / systèmes distribués}\par\vspace{1.2cm}
  {\Huge\bfseries\color{UrbanRed} UrbanOps}\par\vspace{0.35cm}
  {\LARGE\bfseries Architecture multi-middleware et qualité logicielle}\par\vspace{0.2cm}
  {\large Marrakech — signalements, cartographie, alertes, pilotage}\par\vspace{1cm}
  {\color{UrbanRed}\rule{0.5\textwidth}{0.5pt}}\par\vspace{0.5cm}
  {\LARGE\bfseries Rapport professionnel}\par\vspace{0.45cm}
  {\color{UrbanRed}\rule{0.5\textwidth}{0.5pt}}\par\vspace{1.2cm}
  \begin{minipage}{0.88\textwidth}\centering\color{UrbanDark}\large
    Document couvrant la \textbf{méthodologie}, l'\textbf{analyse des besoins}, la \textbf{conception UML}, l'\textbf{implémentation} (Spring Boot~3, Next.js~14), les \textbf{quatre middleware} (REST, JMS, RMI, SOAP), la \textbf{sécurité JWT}, l'\textbf{IA} (Gemini + repli), les \textbf{tests JUnit~5}, \textbf{JaCoCo} et \textbf{SonarQube}.
  \end{minipage}\par\vfill
  {\color{UrbanGreen}\rule{0.78\textwidth}{0.55pt}}\par\vspace{0.55cm}
  \renewcommand{\arraystretch}{1.32}
  \begin{tabular}{@{}>{\bfseries\color{UrbanRed}}r @{\hspace{1em}} p{0.5\textwidth}@{}}
    Auteur(s) : & Prénom \textsc{Nom} \\
    Encadrant(e) : & Prénom \textsc{Nom} \\
    Période : & 2025 -- 2026 \\
  \end{tabular}\par
  \vspace{0.75cm}
  {\large\textbf{Version}~2.0 \quad\textemdash\quad \today}
  \vspace*{0.45cm}
  \restoregeometry
\end{titlepage}

\clearpage
\frontmatter

\chapter*{Dédicaces}
\addcontentsline{toc}{chapter}{Dédicaces}
\thispagestyle{empty}
\vspace*{2.5cm}
\begin{center}
\begin{minipage}{0.82\textwidth}\centering\itshape\large
À nos proches et à toutes les personnes qui œuvrent pour des villes plus sûres et mieux informées.\\[1em]
À nos encadrants pour leur exigence et leur soutien.
\end{minipage}
\end{center}
\clearpage

\chapter*{Remerciements}
\addcontentsline{toc}{chapter}{Remerciements}
\begin{center}
\begin{minipage}{0.88\textwidth}
\justifying
Nous remercions chaleureusement notre encadrant(e) pour le temps accordé, la qualité des orientations et la relecture de ce document. Nous adressons nos remerciements au jury pour l'attention portée à notre travail. Les communautés \textbf{Spring}, \textbf{Next.js}, \textbf{PostgreSQL}, \textbf{ActiveMQ Artemis} et \textbf{Google Gemini} ont fourni documentation, exemples et outils open-source indispensables à la réalisation d'UrbanOps.
\end{minipage}
\end{center}
\clearpage

\chapter*{R\'esum\'e}
\addcontentsline{toc}{chapter}{R\'esum\'e}
\thispagestyle{plain}
\begin{resumebox}
\justifying
UrbanOps est une plateforme \textbf{full stack} pour le \textbf{signalement} et le \textbf{suivi d'incidents urbains} (voirie, eau, éclairage, déchets, sécurité). Le backend \textbf{Spring Boot~3} expose une API \textbf{REST} JSON sous \texttt{/api/v1}, sécurisée par \textbf{JWT}, avec persistance \textbf{PostgreSQL} et \textbf{JPA/Hibernate}. Des canaux complémentaires illustrent les \textbf{systèmes distribués} : messagerie \textbf{JMS} (files d'alertes), \textbf{RMI} pour un service d'analyse IA distant, et \textbf{SOAP/WSDL} pour un reporting à contrat XML strict. Le module d'\textbf{analyse textuelle} s'appuie sur \textbf{Google Gemini} avec \textbf{repli déterministe}. Le frontend \textbf{Next.js~14} (TypeScript, Tailwind, Leaflet) offre landing, tableau de bord et carte. Le présent rapport détaille la méthodologie, la conception, l'implémentation, les tests \textbf{JUnit~5} / \textbf{Mockito}, la couverture \textbf{JaCoCo} et le pilotage qualité \textbf{SonarQube}.
\end{resumebox}
\vspace{0.75cm}
\begin{center}
\textbf{Mots-cl\'es :} ville intelligente ; Spring Boot ; Next.js ; PostgreSQL ; JWT ; JMS ; RMI ; SOAP ; Gemini ; JUnit~5 ; JaCoCo ; SonarQube.
\end{center}
\clearpage

\chapter*{Abstract}
\addcontentsline{toc}{chapter}{Abstract}
\thispagestyle{plain}
\begin{center}
\begin{minipage}{0.88\textwidth}
\justifying
\textit{UrbanOps is a full-stack smart-city supervision platform for urban incident reporting and monitoring. The backend uses Spring Boot~3, REST/JSON APIs secured with JWT, and JPA over PostgreSQL. The project demonstrates distributed middleware patterns: asynchronous JMS alert queues (ActiveMQ Artemis embedded), Java RMI for remote AI analysis, and SOAP with XSD/WSDL for formal reporting contracts. AI-assisted triage uses the Google Gemini API with a keyword-based fallback. The Next.js~14 frontend provides dashboards and maps. This report presents methodology, UML design, implementation, automated testing with JUnit~5, coverage gates with JaCoCo, and SonarQube quality analysis.}
\end{minipage}
\end{center}
\vspace{0.75cm}
\begin{center}
\textbf{Keywords:} smart city; Spring Boot; Next.js; PostgreSQL; JWT; JMS; RMI; SOAP; Gemini; JUnit~5; JaCoCo; SonarQube.
\end{center}
\clearpage

\chapter*{Liste des abr\'eviations et acronymes}
\addcontentsline{toc}{chapter}{Liste des abr\'eviations et acronymes}
\begin{table}[H]
  \centering
  \caption{Abréviations utilisées dans le rapport.}
  \label{tab:abbr}
  \renewcommand{\arraystretch}{1.18}
  \begin{tabularx}{0.96\textwidth}{@{}l Y@{}}
    \toprule
    \textbf{Sigle} & \textbf{Définition} \\
    \midrule
    AMQP & Advanced Message Queuing Protocol \\
    API & Application Programming Interface \\
    CRUD & Create, Read, Update, Delete \\
    DTO & Data Transfer Object \\
    IA & Intelligence artificielle \\
    JMS & Java Message Service \\
    JPA & Jakarta Persistence API \\
    JSON & JavaScript Object Notation \\
    JWT & JSON Web Token \\
    ORM & Object-Relational Mapping \\
    REST & Representational State Transfer \\
    RMI & Remote Method Invocation \\
    SOAP & Simple Object Access Protocol \\
    TLS & Transport Layer Security \\
    UML & Unified Modeling Language \\
    WSDL & Web Services Description Language \\
    XSD & XML Schema Definition \\
    \bottomrule
  \end{tabularx}
\end{table}
\clearpage

\tableofcontents
\clearpage
\listoffigures
\clearpage
\listoftables
\clearpage
\FloatBarrier

% =============================================================================
\mainmatter

\ChapterRoman{Cadre g\'en\'eral, probl\'ematique et m\'ethodologie}
\label{chap:I}
\begin{center}
\begin{minipage}{0.92\textwidth}
\justifying
UrbanOps répond à la dispersion des signalements citoyens et aux délais de traitement en centralisant les incidents, en assurant la \textbf{traçabilité} (référence \texttt{INC-xxxx}, statuts, horodatage) et en orientant les \textbf{alertes} vers les autorités compétentes.
\end{minipage}
\end{center}

\SecRoman{Méthodologie de conduite du projet}
Le travail a combiné une approche \textbf{itérative incrémentale} (backlog type \textbf{Scrum} : sprints fondations, métier, qualité/finalisation) et un cycle en \textbf{V} pour les phases de validation (spécifications $\leftrightarrow$ tests, conception $\leftrightarrow$ tests d'intégration). Les livrables intermédiaires comprennent : modèles \textbf{PlantUML}, API documentée \textbf{OpenAPI/Swagger}, exécution \textbf{Maven} (\texttt{mvn test}, \texttt{mvn verify}), rapports \textbf{JaCoCo} et tableaux de bord \textbf{SonarQube}.

\begin{table}[htbp]
  \centering
  \caption{Phases méthodologiques et artefacts associés.}
  \label{tab:methodo}
  \begin{tabularx}{0.95\textwidth}{@{}l Y@{}}
    \toprule
    \textbf{Phase} & \textbf{Activités / artefacts} \\
    \midrule
    Cadrage & Problématique, périmètre, risques, planning (Gantt, PERT). \\
    Analyse & User stories, besoins fonctionnels et non fonctionnels. \\
    Conception & UML : cas d'utilisation, classes, composants, séquences. \\
    Réalisation & Backend Spring Boot, frontend Next.js, configuration. \\
    Validation & Tests unitaires JUnit~5, couverture, revue SonarQube. \\
    Clôture & Démo, rapport, pistes (mobile, IoT, CI/CD). \\
    \bottomrule
  \end{tabularx}
\end{table}

\SecRoman{Objectifs}
\begin{itemize}
  \item Déclaration d'incidents \textbf{géolocalisés} avec médias ;
  \item Tableau de bord, carte \textbf{Leaflet}, filtres et statistiques ;
  \item Démonstration pédagogique de \textbf{quatre middleware} interopérables ;
  \item Sécurité \textbf{stateless} par JWT et rôles (\texttt{ADMIN}, \texttt{CITIZEN}).
\end{itemize}

\ChapterRoman{Analyse des besoins et sp\'ecifications}
\label{chap:II}
\begin{longtable}{@{}p{3.2cm} >{\RaggedRight\arraybackslash}p{11cm}@{}}
  \toprule
  \textbf{Fonction} & \textbf{Description} \\
  \midrule
  \endfirsthead
  \toprule
  \textbf{Fonction} & \textbf{Description} \\
  \midrule
  \endhead
  Authentification & Inscription, login, JWT, profil \texttt{/auth/me}. \\
  Incidents & CRUD, multipart (JSON + fichier), statuts, carte allégée. \\
  Statistiques & Agrégats tableau de bord, taux de résolution. \\
  Alertes & Routage e-mail pour sévérités \texttt{HIGH} / \texttt{MEDIUM} (via JMS). \\
  Administration & Utilisateurs, incidents, alertes (rôle \texttt{ADMIN}). \\
  IA & Classification / sévérité ; repli si API indisponible. \\
  \bottomrule
\end{longtable}

\SecRoman{Besoins non fonctionnels}
Performance (pagination, projections carte), sécurité (\texttt{BCrypt}, filtre JWT), maintenabilité (couches, DTO, OpenAPI), observabilité (\textbf{Actuator} \texttt{/health}, journaux fichier), résilience messagerie (queue persistante).

\ChapterRoman{Conception et architecture}
\label{chap:III}

\begin{archibox}[Vue logique multicouche]
\begin{center}
\renewcommand{\arraystretch}{1.22}
\begin{tabular}{|>{\RaggedRight\arraybackslash}p{0.88\textwidth}|}
\hline
\textbf{Présentation} — Next.js~14, TypeScript, Tailwind CSS, Axios, Leaflet, Recharts. \\
\hline
\textbf{API REST} — contrôleurs Spring MVC, préfixe \texttt{/api/v1}. \\
\hline
\textbf{Services} — transactions, règles métier, IA, e-mail, JMS. \\
\hline
\textbf{Persistance} — Spring Data JPA, spécifications, PostgreSQL. \\
\hline
\textbf{Transversal} — sécurité Spring Security, JWT (JJWT), Springdoc. \\
\hline
\end{tabular}
\end{center}
\end{archibox}

\SecRoman{Schémas et diagrammes (fichiers image)}
Les diagrammes suivants doivent être exportés (PlantUML ou outil équivalent) puis enregistrés sous les noms indiqués dans le dossier \texttt{docs/latex/figures/} — voir aussi \texttt{figures/README\_FIGURES.md}.

\urbanfig{fig:gantt}{figures/gantt1.png}{0.96}{Diagramme de Gantt — planification des sprints / jalons.}
\urbanfig{fig:pert}{figures/pert1.png}{0.94}{Diagramme PERT — chemin critique et dépendances entre tâches.}
\urbanfig{fig:usecase}{figures/usecase1.png}{0.93}{Diagramme de cas d'utilisation — acteurs et paquetages fonctionnels.}
\urbanfig{fig:classdom}{figures/class_domaine1.png}{0.94}{Diagramme de classes — domaine métier (entités JPA).}
\urbanfig{fig:classtech}{figures/class_technique1.png}{0.94}{Diagramme de classes — couches techniques backend.}
\urbanfig{fig:comp}{figures/composants1.png}{0.92}{Diagramme de déploiement / composants — modules et dépendances.}
\urbanfig{fig:seq}{figures/sequence_incident1.png}{0.92}{Diagramme de séquence — création d'un incident (REST, services, persistance, JMS).}
\FloatBarrier

\ChapterRoman{Environnement technique et middleware}
\label{chap:IV}

\begin{table}[htbp]
  \centering
  \caption{Stack principale UrbanOps.}
  \label{tab:stack}
  \begin{tabularx}{0.94\textwidth}{@{}l Y@{}}
    \toprule
    \textbf{Couche} & \textbf{Technologies} \\
    \midrule
    Backend & Java~17, Spring Boot~3.2, Web, Data JPA, Security, Validation, Mail, Actuator, PostgreSQL, JJWT, Springdoc OpenAPI, Lombok, Maven. \\
    Frontend & Next.js~14, React, TypeScript, Tailwind CSS, Axios, Leaflet, Recharts. \\
    Qualité & JUnit~5, Mockito, spring-security-test, JaCoCo, SonarScanner Maven. \\
    \bottomrule
  \end{tabularx}
\end{table}

\section*{Les quatre middleware (hors numérotation romaine)}
\addcontentsline{toc}{section}{Les quatre middleware}

\begin{longtable}{@{}p{2.2cm} p{2.4cm} >{\RaggedRight\arraybackslash}p{9.2cm}@{}}
  \toprule
  \textbf{Middleware} & \textbf{Rôle} & \textbf{Réalisation dans UrbanOps} \\
  \midrule
  \endfirsthead
  \toprule
  \textbf{Middleware} & \textbf{Rôle} & \textbf{Réalisation dans UrbanOps} \\
  \midrule
  \endhead
  REST & HTTP/JSON & Contrôleurs \texttt{@RestController}, Swagger UI, JWT Bearer. \\
  JMS & Asynchrone & \texttt{AlertProducer} / \texttt{AlertConsumer}, file \texttt{urbanops.alert.queue}, Artemis embarqué. \\
  RMI & RPC Java & Registre port 1099, \texttt{AIAnalysisRemote}, client démo \texttt{RmiClientDemo}. \\
  SOAP & XML/WSDL & Spring-WS, XSD \texttt{urbanops.xsd}, endpoint statistiques et recherche par référence. \\
  \bottomrule
\end{longtable}

\begin{notebox}
\textbf{Captures d'écran} (noms de fichiers pour insertion automatique dans \texttt{figures/}) : \texttt{swagger\_ui.png}, \texttt{jms\_logs\_terminal.png}, \texttt{rmi\_client\_demo.png}, \texttt{soap\_wsdl\_browser.png}, \texttt{soap\_postman\_response.png}. Les figures correspondantes suivent immédiatement ; le mode opératoire détaillé est en \cref{ann:guidefig} (et dans \texttt{figures/README\_FIGURES.md}).
\end{notebox}

\urbanfig{fig:swagger}{figures/swagger_ui.png}{0.92}{Interface Swagger UI — documentation interactive et tests d'API.}
\urbanfig{fig:jmslog}{figures/jms_logs_terminal.png}{0.92}{Terminal — traces JMS (producteur / consommateur / e-mail).}
\urbanfig{fig:rmi}{figures/rmi_client_demo.png}{0.88}{Terminal — exécution de \texttt{RmiClientDemo} (classification à distance).}
\urbanfig{fig:soapwsdl}{figures/soap_wsdl_browser.png}{0.82}{Navigateur — document WSDL (\texttt{urbanops.wsdl}).}
\urbanfig{fig:soappm}{figures/soap_postman_response.png}{0.88}{Client HTTP — enveloppe SOAP et réponse XML.}
\FloatBarrier

\setcounter{chapter}{0}

\chapter{Introduction}
\label{chap:intro}
Ce chapitre introduit la lecture du corps du rapport (numérotation arabe). Les chapitres~\textbf{I} à~\textbf{IV} ont posé le cadre, les besoins, la conception et l'environnement middleware.

\section{Problématique}
Comment concevoir une application \textbf{modulaire}, \textbf{sécurisée} et \textbf{interopérable} capable d'agréger les signalements, d'orchestrer des alertes asynchrones et d'exposer à la fois des API modernes (REST) et des interfaces legacy (SOAP) ou distribuées (RMI) ?

\section{Plan du document}
\cref{chap:real} décrit le backend ; \cref{chap:front} le frontend ; \cref{chap:secia} la persistance, la sécurité et l'IA ; \cref{chap:tests} la stratégie de tests et la qualité ; \cref{chap:deploy} le déploiement local ; \cref{chap:concl} conclut. L'\cref{ann:guidefig} regroupe le \textbf{dossier des figures}, l'\textbf{emplacement} de chaque fichier dans ce document et les \textbf{procédures de capture}. L'\cref{ann:endpoints} liste des endpoints REST représentatifs.

\chapter{R\'ealisation — Backend, persistance et services}
\label{chap:real}
\section{Couche web (REST)}
Les contrôleurs (\texttt{AuthController}, \texttt{IncidentController}, \texttt{StatsController}, \texttt{AlertController}, etc.) appliquent la validation Bean Validation, la pagination et les codes HTTP sémantiques (\texttt{201 Created}, \texttt{404 Not Found}).

\section{Services et orchestration}
\texttt{IncidentService} orchestre la création : résolution catégorie/secteur, enregistrement fichier, appel \texttt{AIAnalysisService}, persistance, puis envoi asynchrone d'alerte via \texttt{AlertProducer} (JMS).

\section{Accès aux données}
Repositories Spring Data, \texttt{JpaSpecificationExecutor} pour filtres combinés, projections pour la carte afin de limiter le volume JSON.

\chapter{R\'ealisation — Frontend Next.js}
\label{chap:front}
\section{Structure de l'application}
App Router, pages publiques (landing) et protégées (dashboard, carte, formulaires), gestion d'état côté client et routage.

\section{Intégration API}
Client HTTP centralisé, en-tête \texttt{Authorization: Bearer} injecté depuis le stockage local après login ; envoi \texttt{multipart/form-data} pour création d'incident avec pièce jointe.

\begin{lstlisting}[caption={Principe d'attachement du jeton (schéma).},label=lst:axios]
api.interceptors.request.use((config) => {
  const token = localStorage.getItem('urbanops_token');
  if (token) config.headers.Authorization = 'Bearer ' + token;
  return config;
});
\end{lstlisting}

\urbanfig{fig:landing}{figures/landing_page.png}{0.88}{Capture — page d'accueil (landing).}
\urbanfig{fig:dash}{figures/dashboard_capture.png}{0.9}{Capture — tableau de bord administrateur.}
\urbanfig{fig:map}{figures/map_incidents.png}{0.9}{Capture — carte des incidents (Leaflet).}
\FloatBarrier

\chapter{Persistance, s\'ecurit\'e et intelligence artificielle}
\label{chap:secia}
\section{Modèle relationnel}
Hibernate génère ou met à jour le schéma selon la configuration ; index sur statuts et coordonnées pour les requêtes carte et filtres.

\section{Sécurité}
\texttt{SecurityFilterChain} : désactivation CSRF pour API stateless, filtres \texttt{JwtAuthenticationFilter}, encodage mots de passe \texttt{BCryptPasswordEncoder}, autorisation par rôle sur les chemins sensibles.

\section{Module IA}
Appel HTTP structuré vers Gemini (réponse JSON attendue), parsing défensif, \textbf{fallback} par mots-clés si quota, latence ou panne réseau.

\chapter{Tests automatis\'es, couverture et qualit\'e}
\label{chap:tests}
\section{Stratégie de tests}
Les tests \textbf{unitaires} (JUnit~5) isolent les services avec \textbf{Mockito} (\texttt{@ExtendWith(MockitoExtension.class)}, \texttt{@Mock}, \texttt{@InjectMocks}). Les tests de contrôleurs peuvent utiliser \texttt{@WebMvcTest} / \texttt{MockMvc}. Les méthodes dépendant d'API externes peuvent être marquées \texttt{@Disabled} et documentées comme tests d'intégration à exécuter manuellement ou en CI avec secrets.

\section{Périmètre observé dans le dépôt}
Classes présentes (indicatif) : \texttt{IncidentServiceTest}, \texttt{UserServiceTest}, \texttt{StatsServiceTest}, \texttt{AIAnalysisServiceTest}, \texttt{AuthControllerTest}. Des extensions possibles : \texttt{AlertServiceTest}, \texttt{IncidentControllerTest}, tests \texttt{@DataJpaTest} sur les repositories.

\section{JaCoCo}
Le plugin \textbf{JaCoCo} (0.8.11) est relié aux phases Maven ; un objectif de couverture minimale sur les packages applicatifs peut être défini dans \texttt{pom.xml}. Commande typique : \texttt{mvn clean verify} ; rapport HTML sous \texttt{target/site/jacoco/}.

\section{SonarQube}
Le plugin Maven Sonar envoie métriques, duplications et résultats de couverture (chemin XML JaCoCo configuré). Les exclusions de fichiers « boilerplate » (DTO, entités, enums, configuration) sont paramétrables pour concentrer l'analyse sur la logique métier.

\urbanfig{fig:jacoco}{figures/jacoco1.png}{0.95}{Rapport JaCoCo — vue par package ou globale.}
\urbanfig{fig:sonar}{figures/sonarqube1.png}{0.95}{Tableau de bord SonarQube — fiabilité, sécurité, dette, couverture.}
\FloatBarrier

\begin{table}[htbp]
  \centering
  \caption{Indicateurs SonarQube à renseigner après analyse (valeurs mesurées).}
  \label{tab:sonar}
  \begin{tabularx}{0.92\textwidth}{@{}l X@{}}
    \toprule
    \textbf{Indicateur} & \textbf{Valeur / commentaire} \\
    \midrule
    Bugs / fiabilité & \ldots \\
    Vulnérabilités & \ldots \\
    Dette technique & \ldots \\
    Couverture & \ldots \\
    Duplications & \ldots \\
    \bottomrule
  \end{tabularx}
\end{table}

\chapter{D\'eploiement et exploitation locale}
\label{chap:deploy}
\begin{lstlisting}[style=bashurban,caption={Démarrage backend.}]
cd backend
mvn spring-boot:run
\end{lstlisting}
\begin{lstlisting}[style=bashurban,caption={Démarrage frontend.}]
cd frontend
npm install
npm run dev
\end{lstlisting}

\chapter{Conclusion et perspectives}
\label{chap:concl}
UrbanOps illustre une architecture \textbf{3~tiers} enrichie par plusieurs styles d'intégration (REST, messagerie, RPC, SOAP). Les perspectives incluent application mobile, notifications push, conteneurisation (Docker/Kubernetes), pipelines CI/CD et extension des suites de tests pour une conformité académique renforcée.

\clearpage
\chapter*{Bibliographie / Webographie}
\addcontentsline{toc}{chapter}{Bibliographie / Webographie}
\begin{thebibliography}{99}
  \bibitem{spring} Spring — \url{https://spring.io/projects/spring-boot}
  \bibitem{next} Next.js — \url{https://nextjs.org/docs}
  \bibitem{gemini} Google AI Gemini — \url{https://ai.google.dev/docs}
  \bibitem{plantuml} PlantUML — \url{https://plantuml.com/}
  \bibitem{sonar} SonarQube — \url{https://docs.sonarsource.com/}
  \bibitem{jjwt} JJWT — \url{https://github.com/jwtk/jjwt}
\end{thebibliography}

\appendix
\chapter{Figures et captures d'écran — dossier, emplacements, procédures}
\label{ann:guidefig}

\section*{Dossier sur le disque}
\addcontentsline{toc}{section}{Dossier sur le disque (figures)}
Toutes les images (exports PlantUML et captures) doivent être placées dans le même répertoire que le fichier \texttt{rapport\_urbanops.tex}, sous-dossier \texttt{figures/} : chemin type \texttt{urbanops/docs/latex/figures/}. La compilation utilise \texttt{\textbackslash includegraphics\{figures/nom.png\}} : le nom doit correspondre \textbf{exactement} au tableau ci-dessous.

{\small
\begin{longtable}{@{}p{3.05cm} p{3.35cm} >{\RaggedRight\arraybackslash}p{8.65cm}@{}}
  \toprule
  \textbf{Fichier} & \textbf{Où dans ce rapport} & \textbf{Comment l'obtenir} \\
  \midrule
  \endfirsthead
  \toprule
  \textbf{Fichier} & \textbf{Où dans ce rapport} & \textbf{Comment l'obtenir} \\
  \midrule
  \endhead
  \texttt{gantt1.png} & Ch.~III, fig.~\ref{fig:gantt} & Exporter le Gantt depuis \texttt{docs/plantuml/01-gantt-urbanops.puml} (PlantUML). \\
  \texttt{pert1.png} & Ch.~III, fig.~\ref{fig:pert} & Exporter \texttt{05-pert-urbanops.puml}. \\
  \texttt{usecase1.png} & Ch.~III, fig.~\ref{fig:usecase} & Exporter \texttt{04-usecase-urbanops.puml}. \\
  \texttt{class\_domaine1.png} & Ch.~III, fig.~\ref{fig:classdom} & Exporter \texttt{02-class-domain.puml}. \\
  \texttt{class\_technique1.png} & Ch.~III, fig.~\ref{fig:classtech} & Exporter \texttt{03-class-backend-technical.puml}. \\
  \texttt{composants1.png} & Ch.~III, fig.~\ref{fig:comp} & Exporter \texttt{06-component-urbanops.puml}. \\
  \texttt{sequence\_incident1.png} & Ch.~III, fig.~\ref{fig:seq} & Exporter \texttt{07-sequence-create-incident.puml}. \\
  \texttt{swagger\_ui.png} & Ch.~IV, fig.~\ref{fig:swagger} & Backend démarré ; ouvrir \texttt{/api/v1/swagger-ui/index.html} ; login ; \textbf{Authorize} avec \texttt{Bearer <token>} ; capture navigateur. \\
  \texttt{jms\_logs\_terminal.png} & Ch.~IV, fig.~\ref{fig:jmslog} & Après \texttt{POST /incidents} (sévérité HIGH/MEDIUM), capture des logs \texttt{[JMS PRODUCER]} / \texttt{[JMS CONSUMER]}. \\
  \texttt{rmi\_client\_demo.png} & Ch.~IV, fig.~\ref{fig:rmi} & Exécuter le client \texttt{RmiClientDemo} (port 1099) ; capture terminal. \\
  \texttt{soap\_wsdl\_browser.png} & Ch.~IV, fig.~\ref{fig:soapwsdl} & Navigateur sur \texttt{/api/v1/soap/urbanops.wsdl}. \\
  \texttt{soap\_postman\_response.png} & Ch.~IV, fig.~\ref{fig:soappm} & \texttt{POST /api/v1/soap} avec corps XML ; capture requête + réponse. \\
  \texttt{landing\_page.png} & \nameref{chap:front}, fig.~\ref{fig:landing} & Capture page d'accueil Next.js (navigateur). \\
  \texttt{dashboard\_capture.png} & \nameref{chap:front}, fig.~\ref{fig:dash} & Capture tableau de bord connecté admin. \\
  \texttt{map\_incidents.png} & \nameref{chap:front}, fig.~\ref{fig:map} & Capture carte Leaflet avec marqueurs. \\
  \texttt{jacoco1.png} & \nameref{chap:tests}, fig.~\ref{fig:jacoco} & Ouvrir \texttt{backend/target/site/jacoco/index.html} après \texttt{mvn verify}. \\
  \texttt{sonarqube1.png} & \nameref{chap:tests}, fig.~\ref{fig:sonar} & Capture tableau de bord SonarQube du projet. \\
  \bottomrule
\end{longtable}
}

\chapter{Endpoints REST repr\'esentatifs}
\label{ann:endpoints}
\begin{longtable}{@{}l l >{\RaggedRight\arraybackslash}p{7.6cm}@{}}
  \toprule
  \textbf{Méthode} & \textbf{Ressource} & \textbf{Description} \\
  \midrule
  \endfirsthead
  \toprule
  \textbf{Méthode} & \textbf{Ressource} & \textbf{Description} \\
  \midrule
  \endhead
  POST & \texttt{/auth/login} & Obtention JWT \\
  POST & \texttt{/auth/register} & Inscription citoyen \\
  GET & \texttt{/incidents} & Liste paginée + filtres \\
  POST & \texttt{/incidents} & Création multipart \\
  GET & \texttt{/incidents/map} & Projection légère carte \\
  GET & \texttt{/stats/dashboard} & Agrégats \\
  GET & \texttt{/stats/jms/status} & Statut broker JMS \\
  GET & \texttt{/stats/rmi/status} & Statut service RMI \\
  GET & \texttt{/soap/urbanops.wsdl} & Contrat WSDL \\
  \bottomrule
\end{longtable}

\end{document}
